%%=============================================================================
%% Conclusie
%%=============================================================================

\chapter{Conclusie}%
\label{ch:conclusie}

% TODO: Trek een duidelijke conclusie, in de vorm van een antwoord op de
% onderzoeksvra(a)g(en). Wat was jouw bijdrage aan het onderzoeksdomein en
% hoe biedt dit meerwaarde aan het vakgebied/doelgroep? 
% Reflecteer kritisch over het resultaat. In Engelse teksten wordt deze sectie
% ``Discussion'' genoemd. Had je deze uitkomst verwacht? Zijn er zaken die nog
% niet duidelijk zijn?
% Heeft het onderzoek geleid tot nieuwe vragen die uitnodigen tot verder 
%onderzoek?

Deze bachelorproef vertrok vanuit een concrete nood binnen de Vlaamse eerstelijnszorg: 
de bestaande instrumenten voor de toegankelijkheidsevaluatie van huisartsenpraktijken worden door huisartsen en praktijkmedewerkers als te tijdrovend, complex en inconsistent ervaren. 
De centrale onderzoeksvraag luidde: \textit{In welke mate kan een proof-of-concept interface bestaande toegankelijkheidsevaluaties voor huisartsen en praktijkmedewerkers sneller, consistenter en gebruiksvriendelijker maken dan de huidige vragenlijsten?}
\medskip

Op basis van de gebruikerstesten met twaalf deelnemers kan worden besloten dat de ontwikkelde proof-of-concept interface een duidelijke meerwaarde biedt. 
De gemiddelde SUS-score van 73,96 situeert de applicatie als \textit{good usable} volgens de internationale normen van Brooke (1995). 
Hoewel deze score iets onder de initieel vooropgestelde drempel van 75 ligt, bevestigt zij dat de interface door huisartsen en praktijkmedewerkers zonder technische achtergrond zelfstandig en vlot kan worden ingevuld. 
De lage scores op de negatieve SUS-stellingen, technische ondersteuning nodig (Q4: 1,08) en veel moeten leren voor gebruik (Q10: 1,17), tonen aan dat de instapdrempel laag is wat essentieel is voor een drukbezet doelpubliek.
\medskip

Wat betreft snelheid bevestigen de tijdsmetingen via Microsoft Clarity dat het invullen van de eigenlijke vragenlijst relatief weinig tijd vraagt (mediaan: 2 min 28 sec). 
Opvallend is dat de eenmalige registratiefase meer tijd in beslag neemt (mediaan: 3 min 15 sec). 
Dit wijst erop dat een toekomstige koppeling met bestaande datasystemen zoals RIZIV-registers of mutualiteitsdata de totale doorlooptijd significant kan reduceren en de adoptiedrempel verder kan verlagen.
\medskip

De kwalitatieve feedback bevestigt de inhoudelijke meerwaarde van de interface. 
Deelnemers waardeerden de gestructureerde opbouw, de overzichtelijke navigatie en de visuele ondersteuning bij het invulproces. 
Tegelijk werden concrete verbeterpunten geïdentificeerd, waaronder het ontbreken van automatisch opslaan, inconsistenties in de interface-elementen (toggles versus knoppen) en de nood aan een meer gepersonaliseerde rapportgeneratie.
\medskip

De lage score op de stelling over regelmatig gebruik (Q1: 2,58) reflecteert geen usabilityprobleem, maar een adoptievraagstuk. 
De meerwaarde van de applicatie voor de dagelijkse praktijk moet duidelijker worden gecommuniceerd, en een structurele koppeling met kwaliteitsverbeteringsprocessen binnen de eerstelijnszone is essentieel voor een duurzame implementatie.
\medskip

De bijdrage van dit onderzoek aan het vakgebied is tweeledig. 
Enerzijds toont het aan dat een gerichte digitale interface, opgebouwd op basis van gevalideerde indicatorensets zoals die van Merckx (2023), het evaluatieproces concreet kan vereenvoudigen voor een doelgroep die structureel onder tijdsdruk staat. 
Anderzijds biedt het onderzoek een methodologisch kader, de combinatie van SUS, Microsoft Clarity-heatmaps en kwalitatieve bevraging, dat als template kan dienen voor toekomstige evaluaties van gelijkaardige zorgtoepassingen.
\medskip

Er zijn echter ook beperkingen. De testgroep bestond uit slechts twaalf deelnemers, allemaal werkzaam binnen Eerstelijnszone Aalst en omliggende regio's. 
Dit beperkt de generaliseerbaarheid van de resultaten naar andere eerstelijnszones of praktijkvormen. 
Een vergelijkende meting met de klassieke papieren of digitale vragenlijst kon bovendien niet worden uitgevoerd, waardoor geen directe tijdsvergelijking mogelijk was. 
Dit vormt een methodologische beperking die in toekomstig onderzoek moet worden opgevangen.
\medskip

Dit onderzoek leidt ook tot nieuwe onderzoeksvragen die verder onderzoek uitnodigen. 
Een eerste vraag betreft de schaalbaarheid: in welke mate is de interface toepasbaar binnen andere eerstelijnszones in Vlaanderen, en welke aanpassingen zijn daarvoor nodig? 
Een tweede vraag sluit aan bij de aanbeveling rond LLM-gebaseerde rapportgeneratie: in welke mate kan een taalmodel zoals de Anthropic API de gegenereerde rapporten relevanter en actiegerichter maken op basis van de ingevulde praktijkgegevens? 
Ten slotte roept de lage score op herhaald gebruik de vraag op hoe de interface beter ingebed kan worden in bestaande kwaliteitszorgprocessen, zodat huisartsen er ook op lange termijn baat bij hebben.
\medskip

Samenvattend kan worden gesteld dat de proof-of-concept interface een veelbelovende eerste stap vormt naar een efficiënter en gebruiksvriendelijker evaluatie-instrument voor de toegankelijkheid van Vlaamse huisartsenpraktijken. 
De resultaten bieden een solide basis voor verdere ontwikkeling en implementatie, en tonen aan dat digitale ondersteuning een reële bijdrage kan leveren aan de kwaliteit van de eerstelijnszorg in Vlaanderen.
